\newpage
\section{Conclusion and Future Work}
In this paper, we proposed a robotic decision-making framework that integrates symbolic knowledge and uncertainty, with the goal of reducing user workload. The proposed method maintains a belief over possible worlds within a POMDP-based planning framework and quantitatively models uncertainty through a probability distribution over these worlds. Based on this, the robot can assess its own uncertainty and selectively query the human only when necessary. This approach enables stable task execution without continuous supervision and minimizes unnecessary interactions, thereby reducing user intervention.

In particular, we leverage entropy-based confidence to determine when to issue queries and select the most informative hypothesis to construct questions, improving the efficiency of interaction. Experimental results show that the proposed method reduces the number of human interactions while maintaining or improving task success rates, and effectively lowers user workload. User studies further support this trend, demonstrating that a system which actively manages uncertainty and requests assistance only when needed can operate effectively in real-world settings.

For future work, several extensions can be considered. First, since the current query mechanism is based on symbolic hypotheses, it can be extended to a more intuitive natural language interface to improve user experience. In addition, incorporating more refined models of perception uncertainty would improve the accuracy of uncertainty estimation. Furthermore, this approach can be extended to multi-robot settings, where multiple agents share uncertainty and coordinate their queries collaboratively.

